% Part: first-order-logic
% Chapter: semantics
% Section: intro-semantics

\documentclass[../../../include/open-logic-section]{subfiles}

\begin{document}

\olfileid{fol}{syn}{its}

\olsection{تعارف}

اظہاراں دے معنی متعین کرنا معنویات دا کم اے۔ معنویات دا مرکزی
تصور !!a{structure} وچ پورا اترن دا اے۔ !!^a{structure} زبان دے
بنیادی حصیاں نوں معنی دیندی اے: !!a{domain} اشیاء دا اک غیر خالی
سیٹ اے۔ مسوّراں دی تعبیر ایس دائرے اُتے چلن دے طور تے ہوندی اے،
!!{constant}s نوں دائرے دے عنصر دتے جاندے نیں، !!{function}s نوں
اوہناں دی مقامیت مطابق !!{domain} دے عناصر دے ترتیب وار ٹپل توں
ایس دائرے دے عناصر ول جان والے فنکشن دتے جاندے نیں، تے
!!{predicate}s نوں !!{domain} اُتے تعلق دتے جاندے نیں۔
% PNB-SEM-001: ماخذ دا دائرے توں اوہدے اپنے آپ ول والا بیان صرف اک-مقامی فنکشن لئی اے؛ ایتھے مقامیت مطابق ٹپل توں دائرے ول فنکشن دتے نیں۔
!!{domain} تے بنیادی لفظیات نوں دتیاں قدراں مل کے !!a{structure}
بناندیاں نیں۔ !!^{variable}s، !!{formula}s وچ آ سکدے نیں؛ معنوی
تعبیر لئی سانوں اوہناں نوں وی !!{domain} دے !!{element}s دینے پیندے نیں---
ایہ متغیر قدر بندی اے۔ پورا اترن دا تعلق آخرکار ایہ ساریاں گلاں
جوڑدا اے۔ !!^a{formula}، !!a{structure}~$\Struct{M}$ وچ
!!a{variable} قدر بندی~$s$ دے لحاظ نال پورا اتر سکدا اے؛
اس نوں $\Sat{M}{!A}[s]$ لکھدے آں۔ ایہ تعلق وی~$!A$ دی ساخت
اُتے استقرا نال تعریف ہوندا اے: منطقی رابطیاں دے صدقی جدول
ورت کے، مثال لئی، $(!A \land !B)$ دا پورا اترنا $!A$ تے~$!B$
دے پورا اترن (یا نہ اترن) دے لحاظ نال تعریف ہوندا اے۔ فیر
ایہ نکلدا اے کہ جے !!{formula}~$!A$، !!a{sentence} ہووے،
یعنی اوہدے وچ کوئی آزاد متغیر نہ ہووے، تاں !!{variable} قدر بندی
دا فرق نہیں پیندا؛ ایس لئی اسیں !!{sentence}s دے !!{structure}s
وچ بس پورا اترن (یا نہ اترن) دی گل کر سکدے آں۔

!!{sentence}s لئی پورا اترن دے تعلق $\Sat{M}{!A}$ دی بنیاد اُتے
اسیں منطقی اعتبار، منطقی نتیجے تے پورا ہون جوگتا دے بنیادی معنوی
تصورات تعریف کر سکدے آں۔ !!^a{sentence} معتبر اے، $\Entails !A$،
جے ہر !!{structure} اوہنوں پورا کرے۔ !!{sentence}s دے کسے سیٹ توں
اوہ منطقی نتیجے دے طور تے نکلدا اے، $\Gamma \Entails !A$، جے ہر
اوہ !!{structure} جیہڑی~$\Gamma$ دے سارے !!{sentence}s پورے کردی
اے،~$!A$ نوں وی پورا کرے۔ تے !!{sentence}s دا کوئی سیٹ پورا ہون
جوگ اے جے کوئی اک !!{structure} اوہدے سارے !!{sentence}s نوں
اکو ویلے پورا کرے۔ کیوں جے !!{formula}s استقرائی طور تے تعریف
ہوندے نیں، تے پورا اترنا وی !!{formula}s دی ساخت اُتے استقرا
نال تعریف ہوندا اے، اسیں استقرا ورت کے اپنی معنویات دیاں خاصیتاں
ثابت کر سکدے آں تے متعین کردہ معنوی تصورات دے باہمی تعلق وکھا
سکدے آں۔

\end{document}
